\label{sec:probes-text}
A probe predicts the difficulty of an item from its features alone, before any label of its benchmark is seen. The target is the mean binary correctness per item, centered within benchmark for fitting. We compare a ridge probe of item embeddings, a probe using log character count, and a sampled random score. We also train two small neural probes (MLPs) on the item embeddings, one on item pass rate and one on the fitted Rasch item difficulty, using 8,462 items from the 11 training benchmarks only.

We score a probe on each benchmark $d$ by two numbers, and average each over benchmarks. Let $y_j$ be the observed difficulty of item $j$ and $\hat y_j$ the probe's prediction. The \emph{Spearman correlation} $\rho_d=\operatorname{corr}\big(\operatorname{rank}(\hat y_j),\operatorname{rank}(y_j)\big)$ compares the two orderings of the items: it is $1$ if the probe orders them perfectly, $0$ if its order is unrelated to the truth, and $-1$ if it is reversed. The \emph{pairwise accuracy} $A_d=\frac{1}{|P_d|}\sum_{(j,k)\in P_d}s_{jk}$ is the share of item pairs the probe orders correctly, where $P_d$ is the set of pairs with $y_j\neq y_k$ (observed ties are excluded), $s_{jk}=1$ if $\hat y_j-\hat y_k$ has the same sign as $y_j-y_k$, $s_{jk}=\tfrac12$ if the probe predicts a tie, and $s_{jk}=0$ otherwise. A random score has $A_d=0.5$ on average. The ID diagnostic holds out items within the 11 training benchmarks; the OOD columns evaluate entire unseen benchmarks. All methods use matched items and sampled pairs.

To test whether a probe helps prediction, we give the item deviations of Rasch a prior mean taken from it,
\begin{equation}\label{eq:probe-prior}
\delta_j\sim\mathcal N\big(\gamma\,z_j,\ \sigma_\delta^2\big),
\end{equation}
where $z_j$ is the probe's predicted difficulty of item $j$, centred within its benchmark (the benchmark rate $\mu_d$ absorbs the average), and $\gamma$ is a scale; $\gamma=0$ is the fitted Rasch. The probes use item features only, never outcomes.

